% !TEX root = ../../hw07.tex
% Source: author's supplied text in chat, 2026-11-06, Q8 (i)--(ii).
% Preserve the logarithmic transformations and differentiated quotients.
% Replace the discrete index by a real variable only in L'Hopital calculations.
\begin{proof}[Solution]
  Part (i).
  Fix $x\in\R$. For all sufficiently large $n$,
  $1+x/n>0$, so the logarithm below is defined. We have
  \[
    \begin{aligned}
      \lim_{n\to\infty}\left(1+\frac xn\right)^n
        &=\lim_{n\to\infty}
          e^{\ln\left((1+\frac xn)^n\right)}\\
        &=\lim_{n\to\infty}e^{n\ln(1+\frac xn)}.
    \end{aligned}
  \]
  Since the exponential function is continuous,
  once the exponent has a finite limit, we obtain
  \[
    \lim_{n\to\infty}e^{n\ln(1+\frac xn)}
      =e^{\lim_{n\to\infty}n\ln(1+\frac xn)}.
  \]

  Strictly, L'H\^opital's rule is applied to the
  corresponding real-variable function, whose limit also determines
  the limit along integer $n$. Write $t$ for this real variable.
  The quotients below are of $0/0$ form as $t\to+\infty$.
  To use the finite-endpoint rule, set $s=1/t\to0^+$. The derivative
  ratio with respect to $s$ equals the displayed ratio with respect
  to $t$, because the common chain-rule factor $-1/s^2$ cancels.

  Solving for the limit of the exponent gives
  \[
    \begin{aligned}
      \lim_{t\to+\infty}t\ln\left(1+\frac xt\right)
        &=\lim_{t\to+\infty}
          \frac{\ln(1+\frac xt)}{1/t}\\
        &\overset{\mathrm{L'H}}{=}\lim_{t\to+\infty}
          \frac{\frac{-x/t^2}{1+x/t}}{-1/t^2}\\
        &=\lim_{t\to+\infty}\frac{x}{1+x/t}\\
        &=x.
    \end{aligned}
  \]
  The denominator derivative $-1/t^2$ is nonzero,
  and the derivative ratio has the finite limit $x$. Restricting to
  $t=n$ therefore gives the required exponent limit. Hence
  \[
    e^{\lim_{n\to\infty}n\ln(1+\frac xn)}
      =\lim_{n\to\infty}\left(1+\frac xn\right)^n=e^x.
  \]

  For part (ii), we begin by expressing the power using the natural
  logarithm and exponential function.
  The base is positive for all sufficiently large
  $n$. Since the exponential function is continuous everywhere on
  $\R$, we may interchange it with the limit, provided the inner
  limit exists and is finite,
  \[
    \begin{aligned}
      \lim_{n\to\infty}\left(1+\sin\frac\pi n\right)^n
        &=\lim_{n\to\infty}
          e^{\ln\left((1+\sin\frac\pi n)^n\right)}\\
        &=\lim_{n\to\infty}e^{n\ln(1+\sin\frac\pi n)}\\
        &=e^{\lim_{n\to\infty}n\ln(1+\sin\frac\pi n)}.
    \end{aligned}
  \]
  We now evaluate the exponent. Rewriting the product as a quotient
  gives $0/0$ form, allowing us to apply L'H\^opital's rule.
  \[
    \begin{aligned}
      \lim_{t\to+\infty}t\ln\left(1+\sin\frac\pi t\right)
        &=\lim_{t\to+\infty}
          \frac{\ln(1+\sin\frac\pi t)}{1/t}\\
        &\overset{\mathrm{L'H}}{=}\lim_{t\to+\infty}
          \frac{\frac1{1+\sin(\pi/t)}
            \cos(\frac\pi t)(-\frac\pi{t^2})}{-1/t^2}\\
        &=\lim_{t\to+\infty}
          \frac{\pi\cos(\pi/t)}{1+\sin(\pi/t)}\\
        &=\frac{\pi\cos(0)}{1+\sin(0)}\\
        &=\pi.
    \end{aligned}
  \]
  The denominator derivative is again nonzero,
  and the derivative ratio tends to $\pi$. Taking $t=n$ gives the
  same exponent limit for the sequence.
  Therefore, substituting this result into the original expression,
  \[
    \begin{aligned}
      \lim_{n\to\infty}\left(1+\sin\frac\pi n\right)^n
        &=e^{\lim_{n\to\infty}n\ln(1+\sin\frac\pi n)}\\
        &=e^\pi.
    \end{aligned}
  \]
\end{proof}
